\documentclass[11pt,a4paper]{article}
\usepackage[margin=1in]{geometry}
\usepackage{amsmath,amssymb,bm}
\usepackage{hyperref}
\usepackage{graphicx}
\usepackage{booktabs}
\usepackage{enumitem}

\newcommand{\vect}[1]{\bm{#1}}
\newcommand{\unit}[1]{\,\mathrm{#1}}
\newcommand{\fm}{\unit{fm}}
\newcommand{\GeV}{\unit{GeV}}
\newcommand{\MeV}{\unit{MeV}}

\title{\texttt{genQE\_3N} --- Three-Nucleon GCF Event Generator\\[4pt]
       \large Physics \& Pipeline Documentation}
\author{GCF with 3N Project}
\date{\today}

\begin{document}
\maketitle
\tableofcontents
\newpage

%======================================================================
\section{Overview}
%======================================================================

\texttt{genQE\_3N} extends the two-nucleon GCF generator to
\textbf{three-nucleon short-range correlations (3N SRC)}.  It generates
events of the form
\begin{equation}
  e + A \;\longrightarrow\; e' + N_{1} + N_{2} + N_{3} + (A{-}3)^{*},
\end{equation}
where the electron scatters off one nucleon of a correlated triplet.
Only $ppn$ triplets are generated.  The default target is $^{12}$C
(residual $^{9}$Be); \texttt{-A}/\texttt{-Z} select any nucleus with
$A \ge 3$, $Z \ge 2$, $N \ge 1$ whose mass and $(A{-}3,\,Z{-}2)$ residual
mass are in the generator's lookup table
(\texttt{QEGenerator\_3N::LookupNuclearMass}).  For $^{3}$He the triplet
is the whole nucleus: there is no residual and no CM motion.

The key conceptual difference from the 2N generator is that
\emph{contact coefficients times universal functions} are replaced by
\textbf{pre-computed 3D density tables} $\rho(\theta_{ab},\,p_{a},\,p_{b})$
built from Faddeev $^{3}$He ground-state wave functions, used as a
stand-in for the 3N SRC wave function in any nucleus.

%======================================================================
\section{From 2N Contacts to 3N Density Matrices}
%======================================================================

\subsection{2N recap}

In the 2N GCF the pair momentum distribution factorizes as
\begin{equation}
  n^{(ij)}(\vect{k}) = C_{ij}^{A}\,|\varphi_{ij}(k)|^{2},
\end{equation}
separating a nucleus-dependent contact $C_{ij}^{A}$ from a universal function
$\varphi_{ij}$.

\subsection{3N generalization}

For three nucleons no analogous simple factorization exists.  Instead, the
full three-body correlation is encoded in a density matrix that depends on
three independent kinematic variables describing the internal configuration
of the triplet:
\begin{equation}
  \rho_{3N}(\theta_{ab},\,p_{a},\,p_{b}),
\end{equation}
where, in the triplet rest frame,
\begin{itemize}[nosep]
  \item $p_{a} = |\vect{p}_{1}|$: momentum of the odd nucleon (in the
        tables, the neutron),
  \item $p_{b} = |\vect{p}_{2}-\vect{p}_{3}|/2$: \emph{half}-relative
        momentum of the other two (the $pp$ pair),
  \item $\theta_{ab}$: angle between $\vect{p}_{a}$ and $\vect{p}_{b}$.
\end{itemize}
The half-relative convention for $p_{b}$ is the one the tables use:
with it the mean kinetic energy of the AV18 table comes out at
$\langle T\rangle = 45.6\MeV$, while reading the same table with the
full relative momentum $\vect{p}_{2}-\vect{p}_{3}$ gives $32\MeV$.
The code calls the two table axes $k_{\text{cm}}$ ($=p_{a}$) and
$k_{\text{rel}}$ ($=p_{b}$).

The tables are plain-text files, loaded at initialization.

\subsection{Data files}

The interaction number \texttt{u} (second positional argument;
\texttt{-u}/\texttt{--wave-function} in \texttt{gcf\_generate}) selects
the table:
\begin{center}
\begin{tabular}{cl}
\toprule
\texttt{u} & File (in \texttt{array/}) \\
\midrule
1 & \texttt{he3\_3Ndensity\_AV18.dnst}       (AV18) \\
2 & \texttt{he3\_3Ndensity\_AV18UIX.dnst}    (AV18 + Urbana IX) \\
3 & \texttt{he3\_3Ndensity\_AV18TM.dnst}     (AV18 + Tucson--Melbourne) \\
4 & \texttt{he3\_3Ndensity\_CDBonn.dnst}     (CD-Bonn) \\
5 & \texttt{he3\_3Ndensity\_CDBonnTM.dnst}   (CD-Bonn + Tucson--Melbourne) \\
6--10 & \texttt{he3\_3Ndensity\_pD\_*.dnst} (the same five, deuteron channel only) \\
\bottomrule
\end{tabular}
\end{center}
The $p$D tables keep only the deuteron channels of the Faddeev wave
function and share the grid, units and Jacobian of $u=1$--5.  Any other
\texttt{u} aborts.  The older \texttt{AV8/N2LO/G3/AV4\_1D.dnst} tables
are still in \texttt{array/} but are no longer loaded.  The directory is
the absolute path CMake bakes in as \texttt{GCF\_3N\_ARRAY\_DIR}; the
environment variable of the same name overrides it at run time.

Each file holds one value per line, $61 \times 101 \times 101 = 622{,}261$
lines, with $\theta_{ab}$ the outermost and $p_{b}$ the innermost index:
\begin{center}
\begin{tabular}{lcl}
\toprule
Variable & Bins & Range \\
\midrule
$\theta_{ab}$ & 61  & $[0^{\circ},\,180^{\circ}]$, $\Delta\theta = 3^{\circ}$ \\
$p_{a}/\hbar c$ & 101 & $[0,\,10]\fm^{-1}$, $\Delta k = 0.1\fm^{-1}$ \\
$p_{b}/\hbar c$ & 101 & $[0,\,10]\fm^{-1}$, $\Delta k = 0.1\fm^{-1}$ \\
\bottomrule
\end{tabular}
\end{center}
Values are trilinearly interpolated at runtime; $p_{a}$ or $p_{b}$ above
$10\fm^{-1}$ returns zero.

\paragraph{Normalization.}
A table stores the raw (Cartesian) density $|\tilde\varphi(p_{a},p_{b},\theta_{ab})|^{2}$
per unit $d^{3}p_{a}\,d^{3}p_{b}$, which is finite at $p_{a}=p_{b}=0$.  With
momenta in GeV/$c$,
\begin{equation}
  8\pi^{2}\int \rho(\theta_{ab},p_{a},p_{b})\,
     p_{a}^{2}\,p_{b}^{2}\,\sin\theta_{ab}\;dp_{a}\,dp_{b}\,d\theta_{ab}
  = \tfrac{1}{3}
\end{equation}
(measured: 0.33308).  The $1/3$ is the isospin assignment: the table
fixes which nucleon is the neutron, and the three equivalent choices sum
to one.  The $8\pi^{2} = 4\pi\cdot 2\pi$ comes from the orientation of
$\vect{p}_{a}$ and the azimuth of $\vect{p}_{b}$ about it.  Powers of
$\hbar c$ are absorbed into the contact $C$.

%======================================================================
\section{Three-Body Kinematics}
%======================================================================

\subsection{Coordinate system}

The three nucleon momenta are parameterized via two Jacobi-like vectors
$\vect{p}_{a}$, $\vect{p}_{b}$ and a three-body center-of-mass momentum
$\vect{P}_{\text{CM}}$.

\paragraph{Internal frame.}
$\vect{p}_{a}$ is placed along $\hat{z}$ and $\vect{p}_{b}$ at polar
angle $\theta_{ab}$ and azimuth $\phi_{baz}$ about it:
\begin{equation}
  \vect{p}_{a} = (0,\,0,\,p_{a}), \qquad
  \vect{p}_{b} = p_{b}\,(\sin\theta_{ab}\cos\phi_{baz},\;\sin\theta_{ab}\sin\phi_{baz},\;\cos\theta_{ab}).
\end{equation}

\paragraph{Rotation to lab.}
Both vectors are then rotated by $\theta_{a}$ about $\hat{y}$ and by
$\phi_{a}$ about $\hat{z}$, so $(\phi_{a},\,\theta_{a},\,\phi_{baz})$
act as three Euler angles, all sampled uniformly.

\paragraph{Individual momenta.}
\begin{equation}
  \vect{v}_{1} = \vect{p}_{a},
  \qquad
  \vect{v}_{2} = -\tfrac{1}{2}\,\vect{p}_{a} + \vect{p}_{b},
  \qquad
  \vect{v}_{3} = -\tfrac{1}{2}\,\vect{p}_{a} - \vect{p}_{b}.
\end{equation}
so $\vect{p}_{b} = (\vect{v}_{2}-\vect{v}_{3})/2$, matching the tables.
The map from $(\vect{P}_{\text{CM}},\vect{p}_{a},\vect{p}_{b})$ to the
three nucleon momenta has unit determinant, so
$d^{3}p_{1}\,d^{3}p_{2}\,d^{3}p_{3} = d^{3}P_{\text{CM}}\,d^{3}p_{a}\,d^{3}p_{b}$.
After Euler rotation, each nucleon receives an additional
$\vect{P}_{\text{CM}}/3$ boost.  The $(A{-}3)$ residual recoils with
$\vect{p}_{A{-}3} = -\vect{P}_{\text{CM}}$.

\subsection{Center-of-mass motion}

\begin{equation}
  P_{\text{CM},\alpha} \sim \mathcal{N}(0,\;\sigma_{\text{CM}}),
  \qquad
  \sigma_{\text{CM}} = 0.15\GeV/c \ \text{(default)}.
\end{equation}
The default is the $^{12}$C 2N-pair literature value.  \texttt{-s} sets
$\sigma_{\text{CM}}$ directly.  In \texttt{genQE\_3N\_FSI} (\texttt{-K})
and \texttt{gcf\_generate} (\texttt{--kF}), giving a Fermi momentum
$k_{F}$ sets $\sigma_{\text{CM}} = \sqrt{3/5}\,k_{F}$ unless
$\sigma_{\text{CM}}$ is also given, the CM width of three uncorrelated
mean-field nucleons.  \texttt{-C 0} (\texttt{--no-cm}) switches the CM
motion off, and it is always off for $A=3$.

%======================================================================
\section{Event Generation Pipeline}
%======================================================================

\subsection{Step 1: Triplet type assignment}

Currently only $ppn$ triplets are generated.  All three nucleons are
initialized as protons, then exactly one --- chosen uniformly at random
among $N_{1}$, $N_{2}$, $N_{3}$ --- is reassigned to a neutron.  Thus
the leading nucleon $N_{1}$ can be either a proton or a neutron (with
probability $2/3$ and $1/3$ respectively).  The weight carries a factor
of~3 to account for the three equivalent assignments.

\subsection{Step 2: Sample three-body CM momentum}

\begin{equation}
  \vect{P}_{\text{CM}} \sim \mathcal{N}^{3}(0,\,\sigma_{\text{CM}}).
\end{equation}
The Gaussian is the CM distribution itself, so it adds no weight factor.

\subsection{Step 3: Sample orientation (Euler angles)}

\begin{equation}
  \phi_{a} \in [-\pi,\,\pi], \qquad
  \cos\theta_{a} \in [-1,\,1], \qquad
  \phi_{baz} \in [-\pi,\,\pi].
\end{equation}
Phase-space weight factor: $\Delta\phi_{a}\,\Delta(\cos\theta_{a})\,\Delta\phi_{baz} = 8\pi^{2}$.

\subsection{Step 4: Sample relative momenta}

\begin{equation}
  p_{a} \in [0,\,5.0]\GeV/c, \qquad
  p_{b} \in [0,\,5.0]\GeV/c, \qquad
  \theta_{ab} \in [10^{-4},\,\pi].
\end{equation}
Here $p_{b}$ is the half-relative momentum of Sec.~2, the variable the
tables are read in.
Phase-space weight factor: $5.0 \times 5.0 \times (\pi - 10^{-4})$.

\subsection{Step 5: Construct individual momenta}

Build $\vect{v}_{1,2,3}$ from the Jacobi decomposition and apply the
Euler rotation.  The event is rejected ($w=0$) unless every internal
momentum satisfies
\begin{equation}
  k_{\min} \le |\vect{v}_{i}| \le 5\GeV/c, \qquad i = 1,2,3,
\end{equation}
with $k_{\min} = 0.25\GeV/c$ by default, or $k_{F}$ when \texttt{-K}
(\texttt{--kF}) is given: the SRC condition $k > k_{F}$ on the initial
state.  The cut acts before the CM boost.  Then add
$\vect{P}_{\text{CM}}/3$ to each nucleon and set
$\vect{p}_{A{-}3} = -\vect{P}_{\text{CM}}$.

\subsection{Step 6: Residual-nucleus energy}

The $(A{-}3)$ residual (for $^{12}$C, $^{9}$Be) is on shell, with mass
from the lookup table plus an excitation energy ($E^{*}=0$ in the
drivers):
\begin{equation}
  m_{A{-}3} = M(A{-}3,\,Z{-}2) + E^{*},
  \qquad
  E_{A{-}3} = \sqrt{|\vect{p}_{A{-}3}|^{2} + m_{A{-}3}^{2}}.
\end{equation}
For $^{3}$He there is no residual and $E_{A{-}3}=0$.

Spectator energies:
\begin{equation}
  E_{2} = \sqrt{|\vect{v}_{2}|^{2} + m_{N}^{2}},
  \qquad
  E_{3} = \sqrt{|\vect{v}_{3}|^{2} + m_{N}^{2}}.
\end{equation}

Leading nucleon energy from conservation (off shell):
\begin{equation}
  E_{1} = m_{A} - E_{A{-}3} - E_{2} - E_{3}.
\end{equation}

\subsection{Step 7: Electron scattering vertex}

The scattered electron direction is sampled uniformly:
\begin{equation}
  \cos\theta_{k} \in [\cos\theta_{k}^{\max},\,\cos\theta_{k}^{\min}],
  \qquad
  \phi_{k} \in [-\pi,\,\pi],
\end{equation}
with $[\theta_{k}^{\min},\theta_{k}^{\max}] = [5^{\circ},\,40^{\circ}]$ in
the drivers (\texttt{-t} in \texttt{genQE\_3N\_FSI}, \texttt{--theta-e}
in \texttt{gcf\_generate}).  Phase-space weight factor:
$\Delta(\cos\theta_{k})\,\Delta\phi_{k}$.

Define
\begin{equation}
  \vect{X} = \vect{v}_{1} + \vect{p}_{\text{beam}},
  \qquad
  Z = m_{A} + E_{\text{beam}} - E_{2} - E_{3} - E_{A{-}3}.
\end{equation}
Requiring the struck nucleon to be on shell after absorbing the photon,
\begin{equation}
  Z - E_{k} = \sqrt{m_{N}^{2} + (\vect{X} - E_{k}\hat{k})^{2}},
\end{equation}
fixes the scattered electron energy:
\begin{equation}
  E_{k} = \frac{Z^{2} - m_{N}^{2} - |\vect{X}|^{2}}
               {2\bigl(Z - \vect{X}\cdot\hat{k}\bigr)}.
\end{equation}
The event is rejected if $E_{k} \le 0$ or $Z - E_{k} \le 0$, or if
$Q^{2}$ falls outside the optional generation window (\texttt{-q} in
\texttt{genQE\_3N\_FSI}, \texttt{--q2} in \texttt{gcf\_generate}; open
by default).

The virtual photon is $q = p_{\text{beam}} - k$ (four-vectors), and the
lead nucleon receives the photon: $p_{\text{lead}} = (E_{1},\vect{v}_{1}) + q$,
so $E_{\text{lead}} = Z - E_{k}$.

\subsection{Step 8: Jacobian}

Integrating the energy-conservation $\delta$-function over $E_{k}$
produces the Jacobian
\begin{equation}
  J_{\delta} = E_{\text{lead}} - \vect{p}_{\text{lead}}\cdot\hat{k},
  \qquad
  w \;\mathrel{*}=\; \left|\frac{E_{\text{lead}}}{J_{\delta}}\right|,
\end{equation}
from $|\partial f/\partial E_{k}| = J_{\delta}/E_{\text{lead}}$, where the
inner derivative of $(\vect{X} - E_{k}\hat{k})^{2}$ is
$-2\,\vect{p}_{\text{lead}}\cdot\hat{k}$.

\subsection{Step 9: Cross section and density weight}

\begin{equation}
  w \;\mathrel{*}=\; \frac{1}{2}\,(2\pi)^{-7}\,C\,t
    \;\times\; \sigma_{eN}(E_{\text{beam}},\,\vect{k},\,\vect{p}_{\text{lead}})
    \;\times\; \rho_{3N}(\theta_{ab},\,p_{a},\,p_{b}),
\end{equation}
where $C = 0.12$ and $t = 2$ are normalization constants.
$\sigma_{eN}$ is the de~Forest cross section of the 2N generator with
Kelly form factors, cc1 by default (\texttt{-c onshell|cc1|cc2}).

\subsection{Step 10: Density evaluation}

The table always has the neutron as nucleon~1 and the $pp$ pair as
nucleons 2,3.  The sampled triplet is therefore relabelled according to
which of $N_{1},N_{2},N_{3}$ is the neutron, $(\vect{v}'_{1},\vect{v}'_{2},\vect{v}'_{3})$,
and the table is read at
\begin{equation}
  p'_{a} = |\vect{v}'_{1}|, \qquad
  p'_{b} = \tfrac{1}{2}|\vect{v}'_{2}-\vect{v}'_{3}|, \qquad
  \theta'_{ab} = \angle(\vect{v}'_{1},\,\vect{v}'_{2}-\vect{v}'_{3}).
\end{equation}
The relabelling is a pointwise re-evaluation of the same density: the
map between the two Jacobi sets has unit determinant, so it carries no
Jacobian.  Because the table is a density per $d^{3}p_{a}\,d^{3}p_{b}$
and the generator samples $p_{a}$, $p_{b}$ and $\theta_{ab}$ uniformly,
the weight needs the measure of the \emph{sampled} variables,
$d^{3}p_{a}\,d^{3}p_{b} = p_{a}^{2}\,p_{b}^{2}\sin\theta_{ab}\;dp_{a}\,dp_{b}\,d\theta_{ab}\,d\Omega_{a}\,d\phi_{baz}$:
\begin{equation}
  \rho_{3N} = p_{a}^{2}\,p_{b}^{2}\,\sin\theta_{ab} \times
         \text{interpolate}\bigl[\rho_{\text{table}}(\theta'_{ab},\,p'_{a}/\hbar c,\,p'_{b}/\hbar c)\bigr].
\end{equation}
The measure uses the unprimed (sampled) variables, the table the primed
ones.

%======================================================================
\section{Weight Summary}
%======================================================================

Collecting all factors:
\begin{equation}
\begin{split}
  w = {}& 3 \times 8\pi^{2} \times 25\,(\pi - 10^{-4})
      \times \Delta(\cos\theta_{k})\,\Delta\phi_{k}
      \times \left|\frac{E_{\text{lead}}}{J_{\delta}}\right|
      \times \frac{C\,t}{2\,(2\pi)^{7}} \\
    & \times \sigma_{eN}
      \times p_{a}^{2}\,p_{b}^{2}\sin\theta_{ab}\;\rho_{\text{table}}(\theta'_{ab},p'_{a},p'_{b})
      \times \Theta(\text{cuts}),
\end{split}
\end{equation}
where $\Theta(\text{cuts})$ stands for the internal-momentum window, the
$E_{k}$ conditions and the $Q^{2}$ window.  Only events with $w>0$ are
written; \texttt{hAttempts} counts every attempt for per-attempt
normalization.

\subsection{Relation to \texttt{3N\_Generator.pdf}}

The code follows the tables where they disagree with the derivation in
\texttt{3N\_Generator.pdf}:
\begin{itemize}[nosep]
  \item Eq.~(6) defines $\vect{p}_{b} = \vect{p}_{2}-\vect{p}_{3}$; the
        tables, and the code, use $(\vect{p}_{2}-\vect{p}_{3})/2$.  With
        the full relative momentum the map from
        $(\vect{p}_{1},\vect{p}_{2},\vect{p}_{3})$ has $|\det| = 1/8$
        instead of 1.
  \item Eq.~(9) is inverted: the measure is
        $p_{a}^{2}\,p_{b}^{2}\sin\theta_{ab}$, not its reciprocal.
  \item Eq.~(25) leaves out $p_{a}^{2}\,p_{b}^{2}\sin\theta_{ab}$, which the
        Cartesian tables require.
  \item Eq.~(18) dots $\vect{p}'_{1}$ into the whole argument
        $\vect{p}_{1}+\vect{p}_{\text{beam}}-\hat{k}$ instead of the inner
        derivative $-\hat{k}$; the code uses
        $J_{\delta} = E_{\text{lead}} - \vect{p}_{\text{lead}}\cdot\hat{k}$.
\end{itemize}

%======================================================================
\section{Alternative Momentum Parameterization}
%======================================================================

For analysis purposes (Dalitz/ternary diagrams, spectral functions), an
alternative parameterization in terms of momentum fractions is available:
\begin{equation}
  f_{i} = \frac{|\vect{p}_{i}|}{|\vect{p}_{1}|+|\vect{p}_{2}|+|\vect{p}_{3}|},
  \qquad
  f_{1} + f_{2} + f_{3} = 1.
\end{equation}
Here $p_{\text{tot}} = p_{1}+p_{2}+p_{3}$ in the triplet rest frame,
where $\vect{p}_{1}+\vect{p}_{2}+\vect{p}_{3}=0$.  Integrating out the
overall orientation,
$d^{3}p_{1}\,d^{3}p_{2} = 8\pi^{2}\,p_{1}\,p_{2}\,p_{3}\;dp_{1}\,dp_{2}\,dp_{3}$,
and $dp_{1}\,dp_{2}\,dp_{3} = p_{\text{tot}}^{2}\;dp_{\text{tot}}\,df_{1}\,df_{2}$,
so the density per $dp_{\text{tot}}\,df_{1}\,df_{2}$ is
\begin{equation}
  \frac{dN}{dp_{\text{tot}}\,df_{1}\,df_{2}}
  \;\propto\; p_{1}\,p_{2}\,p_{3}\,p_{\text{tot}}^{2}\;\rho_{\text{table}},
\end{equation}
up to the constant area factor of the ternary coordinates; it vanishes
at the edges of the ternary plot.
\texttt{QEGenerator\_3N::get\_rho\_ptot\_f1f2f3} currently uses the
reciprocal, $J = (8/\sqrt{3})/(p_{1}\,p_{2}\,p_{3}\,p_{\text{tot}}^{2})$,
which is wrong (issue \#9).  Event generation does not use it (it calls
\texttt{get\_rho}); only the ternary and $\theta$-plot diagnostics
(\texttt{Ternary.cc}, \texttt{analysis/theta\_plot*.cpp}) do.

%======================================================================
\section{Differences from the 2N Generator}
%======================================================================

\begin{center}
\begin{tabular}{lll}
\toprule
Aspect & \texttt{genQE} (2N) & \texttt{genQE\_3N} (3N) \\
\midrule
Correlated cluster & Pair ($N N$) & Triplet ($N N N$) \\
Structure input    & Contacts $C_{ij}\,\varphi^{2}(k)$ & $^{3}$He Faddeev density $\rho(\theta_{ab},p_{a},p_{b})$ \\
Residual system    & $(A{-}2)^{*}$ & $(A{-}3)^{*}$ \\
Phase space dim.   & 6D (CM + rel) & 9D (CM + Euler + Jacobi) \\
Electron vertex    & Quadratic solve for $k_{\perp}$ & Direct solve for $E_{k}$ \\
$\sigma_{\text{CM}}$ & 0.10--0.15 GeV/$c$ & 0.15 GeV/$c$, or $\sqrt{3/5}\,k_{F}$ \\
Target nuclei      & $^{2}$H through $^{208}$Pb & $^{12}$C default; $A\ge3$ in the mass table \\
\bottomrule
\end{tabular}
\end{center}

%======================================================================
\section{Summary of Algorithm}
%======================================================================

For each Monte Carlo trial:
\begin{enumerate}
  \item Assign triplet types ($ppn$); weight $\times\,3$.
  \item Sample $\vect{P}_{\text{CM}}$ (Gaussian, $\sigma = 0.15\GeV/c$ by
        default, or $\sqrt{3/5}\,k_{F}$; none for $A=3$).
  \item Sample Euler angles $(\phi_{a},\theta_{a},\phi_{baz})$ uniformly.
  \item Sample Jacobi magnitudes $(p_{a},p_{b})$ and angle $\theta_{ab}$
        uniformly.
  \item Construct $\vect{v}_{1,2,3}$ via Jacobi decomposition + Euler
        rotation; apply $k_{\min} \le |\vect{v}_{i}| \le 5\GeV/c$; add
        the CM boost.
  \item Compute residual and spectator energies; derive $E_{1}$ from
        conservation.
  \item Sample electron direction $(\theta_{k},\phi_{k})$; solve for $E_{k}$;
        apply the optional $Q^{2}$ window.
  \item Compute Jacobian $J_{\delta} = E_{\text{lead}} - \vect{p}_{\text{lead}}\cdot\hat{k}$.
  \item Evaluate $\sigma_{eN}$ (cc1 by default, Kelly form factors).
  \item Relabel so the neutron is nucleon~1, look up $\rho_{\text{table}}$,
        and multiply by $p_{a}^{2}\,p_{b}^{2}\sin\theta_{ab}$.
  \item Multiply all weight factors; store event if $w > 0$.
\end{enumerate}

\end{document}
